\DocumentMetadata{
  pdfversion=2.0,
  pdfstandard=ua-2,
  lang=en-US,
  tagging=on,
  tagging-setup={math/setup=mathml-SE,tabsorder=structure}
}
\documentclass[11pt,letterpaper]{article}
\usepackage[margin=0.68in,includefoot]{geometry}
\usepackage{newcomputermodern}
\usepackage{amsmath,amscd,mathtools}
\usepackage{tikz,adjustbox}
\providecommand{\square}{\mdlgwhtsquare}
\usepackage[hidelinks]{hyperref}
\hypersetup{
  pdftitle={Logic PhD Exam, August 2006},
  pdfauthor={University of Florida Department of Mathematics},
  pdfsubject={Graduate mathematics examination},
  pdfdisplaydoctitle=true,
  bookmarksopen=true
}
\setcounter{secnumdepth}{0}
\setlength{\parindent}{0pt}
\setlength{\parskip}{0.28em}
\setlength{\emergencystretch}{3em}
\setlength{\leftmargini}{2.15em}
\setlength{\leftmarginii}{2.65em}
\setlength{\leftmarginiii}{3em}
\renewcommand{\labelenumi}{\textbf{\theenumi.}}
\pagestyle{empty}
\raggedbottom
\allowdisplaybreaks
\newenvironment{examproblems}[1]{%
  \begin{enumerate}%
  \setcounter{enumi}{#1}%
  \setlength{\topsep}{0.35em}%
  \setlength{\partopsep}{0pt}%
  \setlength{\itemsep}{0.48em}%
  \setlength{\parsep}{0.18em}%
}{\end{enumerate}}
\newenvironment{examsubparts}{%
  \begin{enumerate}%
  \setlength{\topsep}{0.18em}%
  \setlength{\partopsep}{0pt}%
  \setlength{\itemsep}{0.16em}%
  \setlength{\parsep}{0.1em}%
}{\end{enumerate}}
\newenvironment{examinstructions}{%
  \par\begingroup\itshape
}{%
  \par\endgroup\vspace{0.18em}
}
\makeatletter
\renewcommand\section{\@startsection{section}{1}{0pt}%
  {-1.25ex plus -.25ex minus -.1ex}%
  {0.55ex plus .1ex}%
  {\normalfont\large\bfseries}}
\renewcommand{\@maketitle}{%
  \newpage\null\vskip -1em%
  {\footnotesize\bfseries
    \href{https://www.ufl.edu/}{University of Florida}, \href{https://math.ufl.edu/}{Department of Mathematics}\par}%
  \vskip -0.5em%
  \tagstructbegin{tag=Title}%
    {\LARGE\bfseries\href{https://gma.math.ufl.edu/exams/logic-phd/}{Logic PhD Exam}, August 2006\par}%
  \tagstructend%
  \vskip 0.25em%
  \tagmcbegin{artifact}%
    \hrule height 1pt%
  \tagmcend%
  \vskip 1em%
}
\makeatother
\title{Logic PhD Exam, August 2006}
\author{}
\date{}
\begin{document}
\maketitle
\thispagestyle{empty}
\begin{examinstructions}
Please complete seven problems among the following, including one from each section.
\end{examinstructions}
\section{1. General logic}
\begin{examproblems}{0}
\item
State the incompleteness theorem and explain the role of Gödel numbering in its proof.
\item
Sketch the proof of the compactness theorem for first order logic.
\item
Sketch the proof of Tarski’s undefinability of truth.
\end{examproblems}
\section{2. Model theory}
\begin{examproblems}{3}
\item
Prove that the theory of dense linear orders without endpoints is complete.
\item
Find two models which are elementarily equivalent but not isomorphic.
\item
State and prove Tarski’s criterion for elementary submodels.
\end{examproblems}
\section{3. Set theory}
\begin{examproblems}{6}
\item
Sketch the proof of~\(L\) satisfying CH.
\item
State and prove the infinite Ramsey theorem.
\item
Prove that every stationary subset of \(\omega_1\) can be split into \(\aleph_1\) many pairwise disjoint stationary subsets.
\end{examproblems}
\section{4. Computability}
\begin{examproblems}{9}
\item
Show how to construct a simple recursively enumerable set~\(A\) (meaning that \(\omega-A\) has no infinite r.e. subset).
\item
Define the recursively enumerable set \(K=\{a:\{a\}(a)\mathbin{\downarrow}\}\). Use the \(S_1^1\) function to show that~\(K\) is many-one complete.
\item
Prove that the following are equivalent, for any nonempty \(A\subset\omega\):
\begin{examsubparts}
\item[\textnormal{(a)}]
\(A\) is the domain of some partial recursive function~\(\phi\).
\item[\textnormal{(b)}]
\(A\) is the range of some (total) recursive function~\(f\).
\end{examsubparts}
\end{examproblems}
\end{document}
